% Section 25.5: Thermodynamics: there is no metric at all | Part: 5 | Depends on: part:manifolds, part:integration, sec:lc | Status: done | Original lines: 3529-3558
%!TEX root = ../../../main.tex
\subsection{Thermodynamics: there is no metric at all}\label{sec:phys-thermo}

The state space of a gas (coordinates $V$, $S$, \dots) has no natural inner product: it makes no sense to ``dot'' a pressure with a volume. Yet work and heat are line integrals,
\begin{equation}\label{eq:thermo}
\delta W=-p\,\dd V,\qquad\delta Q=T\,\dd S,\qquad\dd U=\delta Q-p\,\dd V,
\end{equation}
which shows that in thermodynamics the pairing ``covector $\times$ displacement'' is the only possible meaning of work. Here $U$ is a \emph{state function} ($\dd U$ is exact) while $\delta Q$ and $\delta W$ are not closed, which is why heat and work depend on the path.

% Figure | used in parts/05-physics/01-intrinsic-vs-metric/05-thermodynamics.tex
%!TEX root = ../../main.tex
\begin{figure}[ht]
\centering
\begin{tikzpicture}[>=Stealth]
  \draw[->] (0,0)--(4.6,0) node[right] {$V$};
  \draw[->] (0,0)--(0,2.8) node[above] {$p$};
  \draw[blue!70!black,fill=blue!8,rounded corners=10pt] (1.5,0.9) rectangle (3.0,2.0);
  \draw[->,red!75!black,thick] (2.45,1.75) arc (80:-250:0.28);
\end{tikzpicture}
\caption{A thermodynamic cycle in the $(V,p)$ plane. Stokes' theorem identifies $\oint p\,\dd V$ with the enclosed area; the (clockwise) cycle delivers net work to the surroundings. No metric is involved.}
\label{fig:cycle}
\end{figure}


By Stokes (Green) in the $(V,p)$ plane,
\begin{align*}
\oint p\,\dd V&=\iint\dd(p\,\dd V) &&\text{(Stokes)}\\
&=\iint\dd p\wedge\dd V &&(\dd(p\,\dd V)=\dd p\wedge\dd V\text{ since }\dd\dd V=0)\\
&=-\iint\dd V\wedge\dd p &&\text{(antisymmetry)}\\
&=-\text{Area} &&\text{if the loop is counter-clockwise in the $(V,p)$ plane,}
\end{align*}
because $\dd V\wedge\dd p$ is the standard (positive) area form there. So a \emph{clockwise} cycle gives $+$Area (as in Figure~\ref{fig:cycle}): this is the net work delivered by the cycle. Since $\oint\dd U=0$, the first law gives $\oint\delta Q=\oint p\,\dd V$: the heat absorbed over a cycle equals the work delivered. Both are nonzero exactly when the loop encloses a nonzero area (if $\delta Q$ were exact, both would vanish).
Carath\'eodory (1909) \cite{Caratheodory} formulated the second law as a statement about the Pfaffian form $\delta Q$: $1/T$ is an integrating factor, $\delta Q/T=\dd S$.
